\problemname{Squash}

Göran has started playing squash with his friends. He is very dedicated and plays $N$ matches per day.
Göran's goal is, of course, to win as many matches and lose as few matches as possible.
His \emph{satisfaction} after a day of playing squash is equal to the number of matches he won minus the number of matches he lost.

In addition to getting better at squash and winning more matches, Göran has an extra trick to increase
his satisfaction: he claims that he only played \emph{seriously} during $K$ different series of
consecutive matches. Only these matches count in the calculation of his satisfaction. Göran always chooses
these $K$ match series to maximize his satisfaction.

A match series always consists of all matches between the $i$-th and the $j$-th match for some chosen
$i$ and $j$, but Göran can let one (or all!) match series be empty (e.g., maybe he took a very serious
water break instead). The consecutive match series he chooses must not overlap, i.e., each match played
can be part of at most one chosen match series.

Given the match history for a day, what will Göran's satisfaction be if he chooses the match series optimally?


\section*{Input}
The first line of input contains the integer $N$ and $K$ ($1 \leq K \leq N \leq 1\,000$), the number
of matches an the number of serious matches.

The following line contains a string with $N$ characters, the results of the matches that Göran played.
The $i$:th of these characters is \texttt{V} if Göran won the $i$:th match, and \texttt{F} if he lost it.

\section*{Output}
Print an integer: the maximum satisfaction that Göran can achieve.

\section*{Points}
Your solution will be tested on several test case groups.
To get the points for a group, it must pass all the test cases in the group.

\noindent
\begin{tabular}{| l | l | p{12cm} |}
  \hline
  \textbf{Group} & \textbf{Point value} & \textbf{Constraints} \\ \hline
  $1$    & $10$      & $K = 1$ \\ \hline
  $2$    & $20$      & $K = 2$ \\ \hline
  $3$    & $30$      & $K = 3$ \\ \hline
  $4$    & $40$      & No additional constraints. \\ \hline
\end{tabular}
